\rhead{CX201: Digital Logic Design Lab}
\phantomsection 
\section*{Digital Logic Design (CS201)}
\label{course:CX201}

\noindent \small{\hyperref[main:scheme]{$\leftarrow$ Back to Scheme}} \vspace{0.5cm}

\noindent \textbf{Credits:} 2 (0-0-2)

\subsection*{Course Content}
\noindent\textbf{Lab 1: Logic Gate Verification} \
Focuses on familiarization with the digital logic simulation environment (e.g., Logisim, Multisim). Students will construct and simulate all basic logic gates (AND, OR, NOT, NAND, NOR, XOR, XNOR) to verify their respective truth tables. \\

\vspace{0.2cm}

\noindent\textbf{Lab 2: Implementation of Boolean Functions} \
Focuses on translating Boolean algebra into digital circuits. Students will implement logic functions directly from their standard Sum-of-Products (SOP) and Product-of-Sums (POS) forms to understand their structural equivalence. \\

\vspace{0.2cm}

\noindent\textbf{Lab 3: Circuit Minimization with K-Maps} \
Focuses on logic optimization. Students will use Karnaugh Maps (K-Maps) to simplify a 4-variable Boolean function and then implement both the original and the simplified circuits to prove that they are functionally identical but with reduced gate count. \\

\vspace{0.2cm}

\noindent\textbf{Lab 4: Design of Arithmetic Circuits} \
Focuses on core computational blocks. Students will design, implement, and test a Half Adder and a Full Adder. These will then be cascaded to construct a 4-bit Ripple Carry Adder to perform binary addition. \\

\vspace{0.2cm}

\noindent\textbf{Lab 5: Design of Data Selectors and Routers} \
Focuses on MSI combinational logic. Students will design and test a 4-to-1 Multiplexer (MUX) to demonstrate data selection and a 1-to-4 Demultiplexer (DEMUX) to demonstrate data distribution. \\

\vspace{0.2cm}

\noindent\textbf{Lab 6: Design of Encoders and Decoders} \
Focuses on data representation and addressing. Students will implement a Decimal-to-BCD Priority Encoder and a 3-to-8 Line Decoder, which are fundamental components for memory addressing and I/O control. \\

\vspace{0.2cm}

\noindent\textbf{Lab 7: Latches and Flip-Flops} \
Focuses on the fundamentals of sequential memory. Students will build a basic SR Latch using cross-coupled NAND gates and then advance to constructing a clocked D-type Flip-Flop, observing the difference between level-triggered and edge-triggered behavior. \\

\vspace{0.2cm}

\noindent\textbf{Lab 8: Flip-Flop Conversions} \
Focuses on the versatility of sequential elements. Students will implement a JK Flip-Flop and a T Flip-Flop by augmenting a D Flip-Flop with external combinational logic, based on their characteristic equations. \\

\vspace{0.2cm}

\noindent\textbf{Lab 9: Registers and Counters} \
Focuses on sequential MSI components. Students will design and simulate a 4-bit Parallel-In, Serial-Out (PISO) Shift Register and a 3-bit synchronous binary up-counter using JK or T flip-flops. \\

\vspace{0.2cm}

\noindent\textbf{Lab 10: Capstone: Finite State Machine (FSM) Design} \
Focuses on the integration of all concepts to create a control unit. Students will design, implement, and test a complete Moore or Mealy type Finite State Machine to detect a specific binary sequence (e.g., "1011") from a serial input.

\newpage
\rhead{Curriculum Schema}